\subsection{What is formalised}\label{sec:doubling_formal}

This appendix is an account of a Lean~4 development, checked against Mathlib
\texttt{v4.31.0}. Every statement below is a theorem of that development; every proof sketch
names the results that the corresponding proof term actually invokes, and nothing else. The
development contains no unproved goal, and each result stated here has been verified to depend
on no axiom beyond Lean's own \texttt{propext}, \texttt{Classical.choice} and
\texttt{Quot.sound}.

\medskip\noindent
\emph{The object.} The graph is \extref{the counter-example graph of the main text}: a source
$s_0$, a ladder $1,2,3,\dots$, a sink $s_f$, and the forward edges $s_0\rightarrow1$,
$j\rightarrow j+1$, $2j\rightarrow j$, $j\rightarrow s_f$. The predecessors of a ladder state
$j$ are $j-1$ and $2j$, so the backward chain either \emph{decrements} or \emph{doubles}, and
the loop closure adds the wrap $s_f\rightarrow s_0$. The backward policy is
$\BackwardPolicy(j\rightarrow2j)=\varepsilon(j)$ and
$\BackwardPolicy(j\rightarrow j-1)=1-\varepsilon(j)$, and the family studied throughout is
\begin{equation}\label{eq:doubling_family}
    \varepsilon_{c,s}(j)\;:=\;\frac{c}{(j+1)^{s}},\qquad s\geq0,\quad 0<c<2^{s} .
\end{equation}
The target row $\widehat\BackwardPolicy(s_f\rightarrow\cdot)$ is a probability supported in
$\{1,\dots,d\}$ of mean $\bar\jmath$. A \emph{truncation at $K$} deletes every doubling edge
$j\rightarrow2j$ with $2j>K$, leaving a chain on $K+2$ states.

\medskip\noindent
\emph{What is not modelled.} The development builds no stochastic process. It works with the
invariant measure through the balance identity across a cut, with the expected hitting time
through the supersolution that dominates it, and with the descent through a recursion on the
state rather than a chain of random variables. Everything below is therefore a statement about
those objects, and \S\ref{sec:doubling_differences} says where that costs something.
